\section{Reducing a curve to a prism}\label{sec:prism}
For the planar input, take Keller's universal quadrilateral \cite[Thm.~4.1]{keller}. Define
\[
  n(t)=\frac{(1-t^2,2t)}{1+t^2},\qquad
  C=\{q\in\mathbb R^2:n(t_i)\cdot q\le d_i,\ 0\le i\le 3\},
\]
where
\[
 (t_i)=(-631/58,-1,7/16,37/56),\qquad
 (d_i)=(0,0,130783/250000,446199/10^6).
\]
Keller proves that every rectifiable plane path of length one admits an
isometric placement in $C$. Thus $pC$ covers every plane path of length at
most $p$, by extending a shorter path and scaling. The theorem permits planar reflections. Keller supplies a Lean proof of this statement for continuous rectifiable paths, including paths with repeated points \cite{kformal}. The archive includes
the upstream proof sources and finite certificates in \texttt{upstream\_keller/}.
Their file hashes and certificate manifest were checked against the pinned
revision. The upstream Lean kernel replay is available separately using
the commands in that directory.

\paragraph{A prism for each curve.}
For a nonplanar continuous curve $\gamma:[0,1]\to\mathbb R^3$, a unit
vector $u$ exists for which $u\cdot\gamma$ attains its two extrema at
four alternating parameter values. This is Ghomi's alternating support
theorem \cite[Cor.~3.2]{ghomi}. Here is a proof of the instance used here. For $k=1,2,\ldots$,
put $f_u(t)=u\cdot\gamma(t)$ and take probability measures on $[0,1]$
with densities proportional to $\exp(k f_u(t))$ and $\exp(-k f_u(t))$,
respectively. The differences of their first and second moments define
a continuous odd map $S^2\to\mathbb R^2$. It must vanish somewhere. Otherwise its
normalization would give an odd map $S^2\to S^1$. On the equator an odd
map to the circle has odd degree (a continuous lift of its argument changes
by an odd multiple of $\pi$ after a half-turn), whereas its extension to the
hemisphere forces degree zero.
Pass to a subsequence of zeros $u_k\to u$ along which both probability
measures converge weakly. The limits have the same moments through degree
two and are supported on the maximum and minimum sets of $f_u$,
respectively. Indeed $f_{u_k}\to f_u$ uniformly. On a compact subset
where $f_u\le\max f_u-3\varepsilon$, the mass for the positive
weight is exponentially small by comparison with an interval of positive
length where $f_u>\max f_u-\varepsilon$; apply this also to $-f_u$.
The two compact extremal sets are disjoint by nonplanarity.
If there were no four alternating times, the ordered union of the sets
would have at most three consecutive runs of signs. The positive
separation between the sets gives gaps between runs. A polynomial with
one simple root in each intervening gap and an appropriate overall sign
would be positive throughout the maximum set and negative throughout the
minimum set. Its degree would be at most two, contradicting equality of
the integrals. This proves the alternation.

Take the $u$-direction as vertical and let $h$ be the difference of the
extreme heights. The variation of the height is at least $3h$. In arclength
coordinates, the integrated vector
$(\|\gamma'_{xy}\|,|\gamma'_z|)$ has norm at most its integrated norm,
which is one. The horizontal projection has length at most
$p=\sqrt{1-9h^2}$. The curve therefore fits into
$pC\times[0,h]$. A planar curve uses $h=0$ instead. A reflection in a
planar placement extends to a proper spatial motion by also reversing the
vertical coordinate. When $p=0$ the projection is a point and the same
containment holds.

